\chapter{Архитектура, EmotionEngine и разговорен систем}
\label{chap:architecture}

\section{Слоевита архитектура}

Целокупниот систем е прикажан на \cref{fig:system-architecture}. Заедничката
граница меѓу симулацијата и hardware не е actuator command, туку валидирана
емоционална состојба. Ова овозможува истиот разговор и EmotionEngine да се
користат во двата режима, без симулациски \code{Twist}, \code{Joy}, hop или
stomp да стигнат до физичкиот робот.

\begin{figure}[htbp]
\centering
\resizebox{\textwidth}{!}{\begin{tikzpicture}[node distance=7mm and 12mm]
  \node[process,text width=31mm] (user) {Корисник\\\code{emotion\_chat.py}};
  \node[process,below=of user,text width=31mm] (chat) {Chat adapter\\turn ordering};
  \node[process,below=of chat,text width=31mm] (engine) {EmotionEngine\\$v$, $a$, category};
  \node[box,below=of engine,text width=31mm] (state) {ROS state 1.1\\валидирана граница};

  \node[process,right=of chat,text width=32mm] (sidecar) {OpenAI sidecar\\Python 3.12};
  \draw[flow] (user) -- (chat);
  \draw[flow] (chat) -- (engine);
  \draw[flow] (engine) -- (state);
  \draw[flow,<->] (chat) -- node[above,font=\scriptsize]{stream/fallback} (sidecar);

  \node[safety,right=22mm of state,text width=32mm] (mapper) {Expression mapper\\bounded intention};
  \node[safety,right=of mapper,text width=32mm] (bridge) {Safety bridge\\arbitration + gates};
  \node[process,right=of bridge,text width=29mm] (gazebo) {Joy controller\\Gazebo Lite3};
  \draw[flow] (state) -- node[above,font=\scriptsize]{simulation} (mapper);
  \draw[flow] (mapper) -- (bridge);
  \draw[flow] (bridge) -- node[above,font=\scriptsize]{Joy} (gazebo);

  \node[physical,below=11mm of mapper,text width=32mm] (uplink) {Hardware uplink\\NDJSON 1.0};
  \node[physical,right=of uplink,text width=32mm] (runner) {Sole MotionSDK owner\\gates + trajectories};
  \node[physical,right=of runner,text width=29mm] (robot) {Physical Lite3\\12 joints};
  \draw[flow] (state.south) |- node[pos=0.72,above,font=\scriptsize]{physical} (uplink.west);
  \draw[flow] (uplink) -- (runner);
  \draw[flow] (runner) -- (robot);
\end{tikzpicture}
}
\caption{Целокупна архитектура и раздвојување на симулацискиот и физичкиот пат.}
\label{fig:system-architecture}
\end{figure}

Архитектурата има пет логички слоеви:

\begin{enumerate}
  \item \textbf{разговорен слој} --- прифаќа input, управува со turn lifecycle,
  streaming, retry, cancellation и fallback;
  \item \textbf{афективен слој} --- EmotionEngine го ажурира $v$, $a$ и $e$;
  \item \textbf{contract слој} --- проверува schema, sequence, turn correlation
  и bounds;
  \item \textbf{expression/safety слој} --- декларативната емоција ја претвора
  во ограничена симулациска или физичка кореографија; и
  \item \textbf{controller слој} --- Lite3 simulator controller или единствениот
  MotionSDK owner ја извршува веќе одобрената намера.
\end{enumerate}

Ниту chat adapter-от ниту EmotionEngine објавуваат joint effort, motor data,
UDP или Joy. Со тоа природнојазичниот дел не е дел од hard real-time
безбедносниот loop.

\section{Headless EmotionEngine}

\code{EmotionEngine} е поддржаниот интеграциски API во EmotionBot. Една
инстанца претставува разговорна сесија и управува со appraisal, state,
personality, memory, local response shaping и опционална seeded randomness.
Default конфигурацијата е deterministic со seed 7 и исклучена randomness.
За тестирање постои експлицитниот облик \code{event:<emotion>}, со кој без
нејасност се активира секоја од деветте категории.

Состојбата е ограничена според \cref{eq:va-bounds} и \cref{eq:emotion-set}.
Malformed input или непозната категорија не создава actuator command. Domain
logic останува во EmotionBot; ROS time, topics, turn IDs и robot behavior
остануваат во \code{emotion\_bot\_ros} и hardware workspace-от.

\section{Разговорен lifecycle и ordering}

Conversation contract 1.0 бара \code{schema\_version}, timestamp, непразен
\code{turn\_id}, позитивен \code{turn\_index} и еден lifecycle event. Текот е
прикажан на \cref{fig:conversation-lifecycle}. User message се appraisal-ира
веднаш по \code{accepted}; assistant message се appraisal-ира само по
\code{completed}. Така partial delta не ја менува емоцијата.

\begin{figure}[htbp]
\centering
\resizebox{\textwidth}{!}{\begin{tikzpicture}[node distance=9mm and 9mm]
  \node[process,text width=20mm] (input) {input\\turn $n$};
  \node[process,right=of input,text width=23mm] (accept) {accepted\\user appraisal};
  \node[process,right=of accept,text width=22mm] (start) {started\\backend};
  \node[process,right=of start,text width=22mm] (delta) {delta$^*$\\stream};
  \node[process,right=of delta,text width=25mm] (terminal) {completed\\assistant appraisal};
  \draw[flow] (input) -- (accept);
  \draw[flow] (accept) -- (start);
  \draw[flow] (start) -- (delta);
  \draw[flow] (delta) -- (terminal);
  \node[danger,below=of delta,text width=28mm] (failure) {retrying\\offline fallback};
  \draw[dashedflow] (start.south) -- (failure.north west);
  \draw[dashedflow] (failure.north east) -- (terminal.south west);

  \node[danger,below=of accept,text width=25mm] (cancel) {cancelled\\superseded};
  \draw[dashedflow] (input.south east) -- (cancel.north west);
  \coordinate (rejectleft) at ($(cancel.south)+(0,-7mm)$);
  \coordinate (rejectright) at (rejectleft -| terminal.south);
  \draw[dashedflow] (cancel.south) -- (rejectleft) -- (rejectright) -- (terminal.south);
  \node[below=1mm,font=\scriptsize,fill=white,inner sep=1.5pt]
    at ($(rejectleft)!0.55!(rejectright)$) {доцната завршница се отфрла};
\end{tikzpicture}
}
\caption{Lifecycle на разговорен turn со retry, fallback и superseding cancellation.}
\label{fig:conversation-lifecycle}
\end{figure}

Coordinator-от чува најмногу шест turns и 6000 characters context, input до
2000 characters, response до 6000 characters и 300 model output tokens. Нов
input прво го означува стариот turn како cancelled. \code{TurnGate} отфрла
late, duplicate и out-of-order completion. Последниот прифатен turn index се
чува на ROS parameter server, па рестартиран chat node продолжува монотоно во
истиот graph.

Овој механизам решава конкретен забележан race: heartbeat на 2 Hz можел да ја
препокрие one-shot assistant appraisal додека interactive client чекал
корелирана состојба. Корекцијата ја задржува matching assistant state за
активниот turn и е покриена со unit и ROS integration test. Ова е пример каде
sequence/turn semantics се подеднакво важни како самата emotion classifier
логика.

\section{Emotion-state contract 1.1}

Состојбата се објавува како compact JSON на \code{/emotion\_bot/state}. Таа
содржи: \code{schema\_version=1.1}, ROS timestamp, nonnegative sequence,
\code{turn\_id}, source, backend, valence, arousal и emotion. Heartbeat ја
задржува истата semantic sequence и turn ID; не се третира како нова емоција.

Пример:

\begin{lstlisting}[language={},caption={Скратен пример на emotion-state 1.1},label={lst:state-contract}]
{"arousal":0.6,"backend":"deterministic","emotion":"joy",
 "schema_version":"1.1","sequence":1,"source":"user",
 "stamp":{"secs":1,"nsecs":938900000},
 "turn_id":"turn-000001","valence":0.7}
\end{lstlisting}

Валидацијата бара валидни bounds и точно една од деветте категории. Малформирана
состојба во симулација дава neutral mapper output; на физичкиот пат invalid или
replayed transport frame не се прифаќа.

\section{OpenAI sidecar и offline fallback}

Live assistant response се генерира преку официјалниот OpenAI SDK и Responses
streaming API. Sidecar-от работи на Python 3.12,
се врзува само на \code{127.0.0.1:8765} и прима bounded JSON на
\code{/v1/stream}. Проектната конфигурација користи \code{gpt-5-mini},
\code{stream=true}, \code{store=false}, minimal reasoning, low verbosity и
20-секунден timeout. SDK retries се исклучени; coordinator-от прави најмногу
еден retry и потоа преминува на deterministic fallback.

API key се пренесува само преку environment variable во disposable container.
Не се запишува во ROS parameter, source, command argument, image или log.
Health endpoint-от објавува само дали key е присутен и дали bridge е во test
mode. Обичните verification targets користат process-boundary test mode и не
прават платен network request.

Физичкиот OpenAI test открил intermittent DNS failure и недоволно долг client
wait window. Потоа sidecar launch-от добил експлицитни resolvers, а client
deadline е изведен од целосниот provider retry window плус fallback margin.
Две live smoke проверки поминале за 8.239 s и 12.015 s пред повторниот физички
разговорен тест. Овој резултат е network/ordering доказ; сам по себе не е доказ
за роботско движење.
